\documentclass[a4paper]{article} 
\usepackage[T1]{fontenc}
\usepackage[utf8]{inputenc} 
%
\usepackage[spanish]{babel}	
\usepackage[fixlanguage]{babelbib}
    \bibliographystyle{babunsrt} %Asigna el lenguaje a español
%
\usepackage{listingsutf8}
\decimalpoint							%Usa como separador decimal el punto en lugar de la coma


\usepackage{amsmath}						%Librería para agregar símbolos matemáticos
\usepackage{amssymb}						%Librería para agregar símbolos matemáticos
\usepackage{graphicx}					%Librería para agregar figuras
\usepackage{subfigure}										%Crea subfiguras

\usepackage{color}						%Paquete empleado para manipular los colores

\usepackage{fancyhdr}											%Encabezados y pies de pagina
\usepackage{flafter}
\usepackage{txfonts}
\usepackage{MnSymbol}
\usepackage{float}												%Obliga a colocar las figuras donde queremos
\usepackage{multicol}	
\usepackage{lscape}
\usepackage{listings}                     %Paquete para agregar código
\usepackage{pdfpages}
\usepackage{hyperref}
%%%%%%%%%%%%%%%%%%%%%%%%%%%%%%%%%%%%%%%%%%%%%%%%%%%%%%%%%%%%% AGREGADOS 
\usepackage{mathtools} 
% \usepackage{lmodern}
\usepackage{multirow} % Necesario para tablas.
\usepackage{multicol} % Necesario para tablas.
\usepackage{enumerate}
\usepackage{rotating}
\usepackage{tikz}
\usepackage{subfig} % Subfiguras.
\newcommand{\grad}{$^{\circ}$}
% \usepackage{parskip}
\usepackage{longtable}
\usepackage[colorinlistoftodos]{todonotes}
\usepackage[colorlinks=true, allcolors=blue]{hyperref}
% \usepackage{datetime}
%%%%%%%%%%%%%%%%%%%%%%%%%%%%%%%%%%%%%%%%%%%%%%%%%%%%%%%%%%%%%
\input{head}
\begin{document}

%-------------------------------
%	TITLE SECTION
%-------------------------------

\fancyhead[C]{}
	\begin{minipage}{5cm}
		\includegraphics[width=6.5cm]{Logo_un.pdf} 
	\end{minipage}
\hrule \medskip % Upper rule
\begin{minipage}{0.295\textwidth} 
\raggedright
\footnotesize
Santiago Jiménez Gómez  \hfill\\   
CC. 1037659981\hfill\\
sajimenezgo@unal.edu.co
\end{minipage}
\begin{minipage}{0.4\textwidth} 
\centering 
\large 
TRABAJO DEL CURSO  \\ 
\normalsize 
Mecánica Estructural, 2020-2\\ 
\end{minipage}
\begin{minipage}{0.295\textwidth} 
\raggedleft
2020-11-6\\
\end{minipage}
\medskip\hrule 
\bigskip

%-------------------------------
%	CONTENTS
%-------------------------------

%---Matrices Locales


%------------------------------------------------

\end{document}

...
